\section{Query \& Retention Policies}
\subsection{Retention Design}
Retention is configured at the bucket level in InfluxDB v2. No retention policy is applied within a bucket; the bucket-level setting is the sole retention mechanism.

\begin{table}[H]
  \caption{Measuremnts}
  \begin{center}
    \begin{tabular}{lll}
      \hline
      \textbf{Bucket} & \textbf{Telemetry} & \textbf{Rationale}\\
      \hline
      \texttt{telemetry} & 90 days & Covers the active thesis validation period at full resolution \\
      \texttt{telemetry\_1m} & Indefinite & Long--term aggregate data for reporting and trend analysis \\
      \hline
    \end{tabular}
  \end{center}
\end{table}

\subsection{Downsampling Task}
A Flux scheduled task aggregates telemetry into telemetry\_1m at 1--minute intervals. mean, min, and max are stored as separate fields to support both trend analysis and power quality event detection within each window.

\begin{lstlisting}[language=flux]
option task = {name: "downsample_1m", every: 1m}

from(bucket: "telemetry")
  |> range(start: -task.lastSuccessTime)
  |> filter(fn: (r) => r._field == "v")
  |> aggregateWindow(every: 1m, fn: mean, createEmpty: false)
  |> set(key: "_field", value: "v_mean")
  |> to(bucket: "telemetry_1m")

from(bucket: "telemetry")
  |> range(start: -task.lastSuccessTime)
  |> filter(fn: (r) => r._field == "v")
  |> aggregateWindow(every: 1m, fn: min, createEmpty: false)
  |> set(key: "_field", value: "v_min")
  |> to(bucket: "telemetry_1m")

from(bucket: "telemetry")
  |> range(start: -task.lastSuccessTime)
  |> filter(fn: (r) => r._field == "v")
  |> aggregateWindow(every: 1m, fn: max, createEmpty: false)
  |> set(key: "_field", value: "v_max")
  |> to(bucket: "telemetry_1m")
\end{lstlisting}

\subsection{Key Query Parameters}
The following Flux examples cover the primary query patterns anticipated during validation.

Time-window data for a single channel:
\begin{lstlisting}[language=flux]
from(bucket: "telemetry")
  |> range(start: -1h)
  |> filter(fn: (r) => r._measurement == "voltage"
                    and r.node_id == "node_1"
                    and r.measurement == "va")
\end{lstlisting}

Cross-node comparison for the same channel:
\begin{lstlisting}[language=flux]
from(bucket: "telemetry")
  |> range(start: -1h)
  |> filter(fn: (r) => r._measurement == "voltage"
                    and r.measurement == "va")
  |> group(columns: ["node_id"])
\end{lstlisting}

Cross--node comparisons are only meaningful where PTP synchronisation is confirmed locked. The ptp\_status measurement should be checked before drawing cross--node conclusions from the historian.

All channels from a single node (debugging):
\begin{lstlisting}[language=flux]
from(bucket: "telemetry")
  |> range(start: -15m)
  |> filter(fn: (r) => r.node_id == "node_1")
\end{lstlisting}


PTP synchronisation quality over a window:
\begin{lstlisting}[language=flux]
from(bucket: "telemetry")
  |> range(start: -1h)
  |> filter(fn: (r) => r._measurement == "ptp_status"
                    and r.node_id == "node_1")
\end{lstlisting}

\subsection{Grafana Integration}
Grafana is a confirmed thesis deliverable. Grafana connects to InfluxDB v2 via the native Flux data source. Key dashboard panels anticipated:

\begin{table}[H]
  \caption{Measuremnts}
  \begin{center}
    \begin{tabular}{lll}
      \hline
      \textbf{Panel} & \textbf{Source Bucket} & \textbf{Query Pattern}\\
      \hline
      Per--phase voltage time--series & telemetry & Single--Channel, cross--node \\
      Per--phase current time--series & telemetry & Single--Channel, cross--node \\
      PTP offset over time (per node) & telemetry & ptp\_status measurement \\
      1--minute voltage mean/min/max & telemetry\_1m & Single--Channel, cross--node \\
      \hline
    \end{tabular}
  \end{center}
\end{table}
Grafana dashboard design is out of scope for this document and will be addressed in a separate deliverable.

